\begin{tikzpicture}[font=\scriptsize, >=latex]
% ELF anatomy of the real ch01_sum.elf
\draw[thick, fill=graybg] (0.6,0.6) rectangle (4.4,4.6);
\node[font=\footnotesize\bfseries] at (2.5,4.3) {ELF header};
\foreach \y/\t/\s in {
  3.8/{.text (0x96 B)}/{the code, vma 0x2011c0},
  3.2/{.data (0x34 B)}/{initialised data},
  2.6/{.eh\_frame}/{unwind info},
  2.0/{.symtab + .strtab}/{names for linkers},
  1.4/{.shstrtab + .comment}/{section names, meta}}{
  \draw[thick] (0.8,\y-0.24) rectangle (4.2,\y+0.24);
  \node[align=left, font=\tiny] at (0.95,\y) {\t};
}
\node[font=\tiny, text=slate] at (2.5,0.9) {real \ttfamily ch01\_sum.elf};
% section headers annotation
\draw[->, thick, slate] (4.4,3.8) -- (6.0,3.8)
  node[right, align=left, font=\tiny, text width=4.2cm]
  {code bytes, executable --- the part your cpu8 would fetch};
\draw[->, thick, slate] (4.4,2.0) -- (6.0,2.0)
  node[right, align=left, font=\tiny, text width=4.2cm]
  {symbol table --- the linker's contract, invisible at run time};
\node[align=left, text=slate, text width=10cm] at (5.3,-0.2)
{An ELF is a table of contents plus the bytes each section points
at: same bytes, four different jobs.};
\end{tikzpicture}
